% from ega1-9-fr.tex
\oldpage[I]{174}

\subsection{Prolongement des faisceaux quasi-cohérents.}
\label{subsection:I.9.4-fr}

\begin{env}[9.4.1]
\label{I.9.4.1-fr}
Soient $X$ un espace topologique, $\mathscr{F}$ un faisceau d'ensembles
(resp. de groupes, d'anneaux) sur $X$, $U$ une partie ouverte de $X$,
$\psi:U\to X$ l'injection canonique, $\mathscr{G}$ un sous-faisceau de
$\mathscr{F}|U=\psi^*(\mathscr{F})$. Comme $\psi_*$ est exact à gauche,
$\psi_*(\mathscr{G})$ est un sous-faisceau de
$\psi_*(\psi^*(\mathscr{F}))$~; si on considère l'homomorphisme canonique
$\rho:\mathscr{F}\to\psi_*(\psi^*(\mathscr{F}))$
(\hyperref[0.3.5.3-fr]{0, 3.5.3}), nous désignerons par
$\overline{\mathscr{G}}$ le sous-faisceau
$\rho^{-1}(\psi_*(\mathscr{G}))$ de $\mathscr{F}$. Il résulte immédiatement
des définitions que pour tout ouvert $V$ de $X$,
$\Gamma(V,\overline{\mathscr{G}})$ est formé des sections
$s\in\Gamma(V,\mathscr{F})$ dont la restriction à $V\cap U$ soit une section
de $\mathscr{G}$ au-dessus de $V\cap U$. On a donc
$\overline{\mathscr{G}}|U=\psi^*(\overline{\mathscr{G}})=\mathscr{G}$, et
$\overline{\mathscr{G}}$ est le \emph{plus grand} sous-faisceau de
$\mathscr{F}$ induisant $\mathscr{G}$ sur $U$~; nous dirons que
$\overline{\mathscr{G}}$ est le \emph{prolongement canonique} du
sous-faisceau $\mathscr{G}$ de $\mathscr{F}|U$ en un sous-faisceau de
$\mathscr{F}$.
\end{env}

\begin{proposition}[9.4.2]
\label{I.9.4.2-fr}
Soient $X$ un préschéma, $U$ une partie ouverte de $X$ telle que l'injection
canonique $j:U\to X$ soit un morphisme quasi-compact
\emph{(ce qui sera vérifié pour \emph{tout} $U$ si l'espace sous-jacent $X$
est \emph{localement noethérien}
(\hyperref[I.6.6.4-fr]{6.6.4, (i)}))}. Alors~:
\begin{enumerate}
\item[(i)] Pour tout $(\mathscr{O}_X|U)$-Module quasi-cohérent
  $\mathscr{G}$, $j_*(\mathscr{G})$ est un $\mathscr{O}_X$-Module
  quasi-cohérent et on a
  $j_*(\mathscr{G})|U=j^*(j_*(\mathscr{G}))=\mathscr{G}$.
\item[(ii)] Pour tout $\mathscr{O}_X$-Module quasi-cohérent $\mathscr{F}$ et
  tout sous-$(\mathscr{O}_X|U)$-Module quasi-cohérent $\mathscr{G}$, le
  prolongement canonique $\overline{\mathscr{G}}$ de $\mathscr{G}$ (9.4.1)
  est un sous-$\mathscr{O}_X$-Module quasi-cohérent de $\mathscr{F}$.
\end{enumerate}
\end{proposition}

\begin{proof}
Si $j=(\psi,\theta)$ ($\psi$ étant l'injection $U\to X$ des espaces
sous-jacents), on a par définition $j_*(\mathscr{G})=\psi_*(\mathscr{G})$
pour tout $(\mathscr{O}_X|U)$-Module $\mathscr{G}$, et en outre
$j^*(\mathscr{H})=\psi^*(\mathscr{H})=\mathscr{H}|U$ pour tout
$\mathscr{O}_X$-Module $\mathscr{H}$ en raison de la définition d'un
préschéma induit sur un ouvert. (i) est donc un cas particulier de
(\hyperref[I.9.2.2-fr]{9.2.2, a)})~; pour la même raison,
$j_*(j^*(\mathscr{F}))$ est quasi-cohérent, et comme
$\overline{\mathscr{G}}$ est l'image réciproque de $j_*(\mathscr{G})$ par
l'homomorphisme $\rho:\mathscr{F}\to j_*(j^*(\mathscr{F}))$, (ii) résulte de
(\hyperref[I.4.1.1-fr]{4.1.1}).
\end{proof}

On notera que l'hypothèse que le morphisme $j:U\to X$ est quasi-compact est
aussi vérifiée lorsque l'ouvert $U$ est \emph{quasi-compact} et $X$ un
\emph{schéma}~: en effet, $U$ est alors réunion finie d'ouverts affines
$U_i$, et pour tout ouvert affine $V$ de $X$, $V\cap U_i$ est un ouvert
affine (\hyperref[I.5.5.6-fr]{5.5.6}), donc quasi-compact.

\begin{corollary}[9.4.3]
\label{I.9.4.3-fr}
Soient $X$ un préschéma, $U$ un ouvert quasi-compact de $X$ tel que le
morphisme d'injection $j:U\to X$ soit quasi-compact. Supposons en outre que
tout $\mathscr{O}_X$-Module quasi-cohérent soit limite inductive de ses
sous-$\mathscr{O}_X$-Modules quasi-cohérents de type fini
\emph{(ce qui a lieu lorsque $X$ est un \emph{schéma affine})}. Soient alors
$\mathscr{F}$ un $\mathscr{O}_X$-Module quasi-cohérent, $\mathscr{G}$ un
sous-$(\mathscr{O}_X|U)$-Module quasi-cohérent de type fini de
$\mathscr{F}|U$. Il existe alors un sous-$\mathscr{O}_X$-Module
quasi-cohérent de type fini $\mathscr{G}'$ de $\mathscr{F}$ tel que
$\mathscr{G}'|U=\mathscr{G}$.
\end{corollary}

\begin{proof}
En effet, on a $\mathscr{G}=\overline{\mathscr{G}}|U$, et
$\overline{\mathscr{G}}$ est quasi-cohérent d'après (9.4.2), donc limite
inductive de ses sous-$\mathscr{O}_X$-Modules quasi-cohérents de type fini
$\mathscr{H}_\lambda$. Par suite $\mathscr{G}$ est limite inductive des
$\mathscr{H}_\lambda|U$, donc égal à un des $\mathscr{H}_\lambda|U$
puisqu'il est de type fini (\hyperref[0.5.2.3-fr]{0, 5.2.3}).
\end{proof}

\begin{remark}[9.4.4]
\label{I.9.4.4-fr}
Supposons que pour \emph{tout} ouvert affine $U\subset X$ le morphisme
d'injection $U\to X$ soit quasi-compact. Alors si la conclusion de (9.4.3) a
lieu pour tout ouvert affine $U$ et tout sous-$(\mathscr{O}_X|U)$-Module
quasi-cohérent de type fini $\mathscr{G}$ de $\mathscr{F}|U$,
